
In this section, we show that a certain fragment of string constraints
whose decidability status has been left open in the literature can be
reduced in logarithmic space to generalised Sem{\"e}nov arithmetic,
and is hence decidable in EXPSPACE. This demonstrates an important
application of our results on generalised Sem{\"e}nov arithmetic, with
deep connections to solving string constraints in practice, which has
been one of the motivations for our work.

Let $\Sigma = \{ 0,1\}$. The \emph{theory of enriched string
  constraints} $\TRElnc$ is the first-order theory of the two-sorted
structure \[\langle \Sigma^*,\N; \{w\}_{w \in \Sigma^*}, \cdot, \slen,
\strnum, \{R_i \}_{i\in \N}, 0,1,+\rangle,\] where
%
\begin{itemize}
\item the binary function $\cdot$ over $\Sigma^*$ is the string
  concatenation operator,
\item the unary function $\slen\colon \Sigma^* \to \N$ returns on
  input $w$ the length $\abs{w}$ of $w$,
\item the unary function $\strnum\colon \Sigma^* \to \N$ on input $u$
  returns $\eval{u}$, and
\item $R_0,R_1,\ldots \subseteq \Sigma^*$ is an enumeration of all
  regular languages.
\end{itemize}
%
The remaining predicates, constant and function symbols are defined in
their standard semantics.

The above theory was introduced in~\cite{BKMMDNG20}, where an SMT solver addressing some fragments of this theory was defined, implemented, and compared to other state of the art solvers which can handle such string constraints. Extending \cite{BKMMDNG20}, \cite{BDGKMMN23} presents in more details the motivation behind considering this theory and its fragments. More precisely, the authors of \cite{BDGKMMN23} analysed an extensive collection of standard real-world benchmarks of string constraints and extracted the functions and predicates occurring in them. The works \cite{BKMMDNG20,BDGKMMN23} focused on benchmarks that do not contain word equations, and the result of the aforementioned benchmark-analysis produced exactly the four functions and predicates mentioned above: $\slen$, $\strnum$, regular language membership, and concatenation of strings. 

Complementing the practical results of \cite{BKMMDNG20}, \cite{BDGKMMN23} showed a series of theoretical results regarding fragments of $\TRElnc$. In particular, the existential theory of $\TRElnc$ is shown to be undecidable. Moreover, \cite{BDGKMMN23} leaves as an open problem the question whether the existential theories of $\TREln$ and $\TREnc$, which drop the concatenation operator and length function, respectively, are decidable. From these two, the existential theories of $\TREln$ seems particularly interesting, as all instances from the benchmarks considered in the analysis \cite{BDGKMMN23} can be easily translated into a formula from this particular fragment of $\TREln$. Indeed, by the results reported in Table 1.b from \cite{BDGKMMN23},  %(which we have also double checked), 
no instance contains both concatenation of strings and the $\strnum$ function; moreover, the concatenation of strings, which appears only in formulas involving regular membership predicates and, in some cases, length function, can be easily removed in all cases by a folklore technique called automata splitting (see, e.g., \cite{AACHRRS15}). Therefore, showing that the existential fragment of $\TREln$ is decidable essentially shows that one can decide all the instances from the standard benchmarks analysed in \cite{BDGKMMN23}.

In this paper, we solve this open problem. By a reduction to generalised Sem{\"e}nov arithmetic, we can settle the decidability status of $\TREln$:
%
\begin{theorem}
  The existential fragment of $\TREln$ is decidable in EXPSPACE.
\end{theorem}
%
%% We equivalently interpret both the unary functions $\slen$
%% and $\strnum$ and the binary function $+$ as binary and ternary
%% relations over $\Sigma^*\times\N$ and $\N^3$, respectively.

Again, we treat $\TREln$ as a relational structure. Without loss of
generality, we may assume that atomic formulas of $\TREln$ are one of
the following:
%
\begin{itemize}
\item $R(s)$ for some string variable $s$ and a regular language $R$;
\item $s=t$ for some string variables $s$ and $t$;
\item $\slen(s,x)$ or $\strnum(s,x)$ for some string variable $s$ and
  integer variable $x$; or
\item $\vec a \cdot \vec x \ge b$ for a vector of integer variables
  $\vec x$.
\end{itemize}
%
The size of a formula of $\TREln$ is defined in the standard way as
the number of symbols required to write it down, assuming binary
encoding of numbers, and where the size of some $R$ is the size of the
smallest DFA accepting $R$. Furthermore, in a quantifier-free formula
$\varphi$ of $\TREln$, we may without loss of generality assume that
all atomic formulas occur positive, except for atomic formulas $s=t$.

We now describe the reduction to existential Sem{\"e}nov
arithmetic. The idea underlying our proof is that we map a string $s$
to the number $\eval{1s}$. Note that we cannot directly treat strings
in $\Sigma^*$ as natural numbers due to the possibility of leading
zeros. This encoding enables us to treat strings as numbers and to
implement the functions $\strnum$ and $\slen$ in generalised
Sem{\"e}nov arithmetic. Given a quantifier-free formula $\varphi$ of
$\TREln$, we define by structural induction on $\varphi$ a function
$\sigma$ that maps $\varphi$ to an equi-satisfiable formula of
generalised Sem{\"e}nov arithmetic:
%
\begin{itemize}
\item Case $\varphi \equiv R(s)$: $\sigma(\varphi)\defeq (0^*1 R)(s)$;
\item Case $\varphi \equiv s=t$ or $\varphi \equiv \neg(s=t)$:
  $\sigma(\varphi) \defeq s=t$ or $\sigma(\varphi) \defeq \neg s=t$,
  respectively;
\item Case $\varphi \equiv \strnum(s,x)$: $\sigma(\varphi)\defeq
  \exists y.\, 2^y\leq s\land s< 2^{y+1} \land x = s-2^y$;
\item Case $\varphi \equiv \slen(s,x)$: $\sigma(\varphi)\defeq 2^{x}\le s\land s < 2^{x+1}$;
\item Case $\varphi \equiv \vec a\cdot \vec x \ge b$: $\sigma(\varphi) \defeq \vec a\cdot \vec x \ge b$; and
\item Case $\varphi \equiv \varphi_1 \sim \varphi_2$, ${}\sim{} \in \{
  \wedge, \vee \}$: $\sigma(\varphi)\defeq \sigma(\varphi_1) \sim
  \sigma(\varphi_2)$.
\end{itemize}
%

\begin{lemma}
  Let $\varphi$ be a quantifier-free formula of $\TREln$ and $S$ be
  the set of string variables occurring in $S$. Then $\varphi$ is
  satisfiable if and only if $\sigma(\varphi) \wedge \bigwedge_{s\in
    S} s>0$ is satisfiable.
\end{lemma}
\begin{proof}
  Observe that the variables occurring in $\varphi$ are the same
  variables as those occurring in $\sigma(\varphi)$. Let $S$ be the
  set of string variables in $\varphi$ and $X$ be the set of
  integer-valued variables in $\varphi$. Given an assignment $\mathcal
  I_S \colon S \to \{0,1\}^*$, we define $\tilde{\mathcal{I}}_S \defeq
  S \to \N$ such that $\tilde{\mathcal I}_S(s) \defeq \eval{1 \mathcal
    I_S(s)}$. Subsequently, denote by $\mathcal I_X\colon X \to \N$ an
  assignment to the integer-valued variables. We show by structural
  induction on $\varphi$ that $(\mathcal I_S,\mathcal I_x) \models
  \varphi$ if and only if $(\tilde{\mathcal I}_S, \mathcal I_X)
  \models \sigma(\varphi) \wedge \bigwedge_{s\in S} s>0$:
%
  \begin{itemize}
  \item Case $\varphi \equiv R(s)$: Let $\mathcal I_S(s)=b_{n-1}\cdots
    b_0$, we have $\mathcal I_S(s)\in R$ if and only if $2^n +
    \sum_{i=0}^{n-1} 2^ib_i \in \eval{0^*1R}$, noting that $2^n +
    \sum_{i=0}^{n-1}2^i b_i=\eval{1b_{n-1}\cdots b_0}=\tilde{\mathcal
      I}_S(s)$.
  \item Case $\varphi \equiv \strnum(s,x)$: Let $\mathcal
    I_S(s)=b_{n-1}\cdots b_0$ and $\mathcal I_X(x)=m$. We have that
    $m=\sum_{i=0}^{n-1}2^ib_i$ if and only if $m=\tilde{\mathcal
      I}_S(s)-2^n$ if and only if $(\tilde{\mathcal I}_S,\mathcal
    I_X)\models \sigma(\varphi) \wedge \bigwedge_{s\in S} s>0$.
  \item Case $\varphi\equiv \slen(s,x)$: Let $\mathcal
    I_S(s)=b_{n-1}\cdots b_0$ and $\mathcal I_X(x)=m$. We have that
    $m=n$ if and only if $2^m\le \eval{1b_{n-1}\cdots b_0} <2^{m+1}$
    if and only if $(\tilde{\mathcal I}_S,\mathcal I_X)\models
    \sigma(\varphi) \wedge \bigwedge_{s\in S} s>0$.
  \end{itemize}
  The remaining cases follow obviously.
\end{proof}
